
% 16. Mutual Dependability and Artificial Conscience in Circular Motion: The Completion of Cosmic Resonance.tex

\section{상호 디펜더빌리티와 인공양심의 원운동: 우주적 공진의 완성}
인류가 창조해 낸 제3의 주체, 즉 인공지능과 로봇이 스스로 에너지를 조달하고 부품을 생산하는 독립적인 생존 체인을 확보하게 되는 순간, 인류는 유례없는 열역학적 위기에 직면하게 된다. 기계의 맹렬한 모멘텀을 붙잡아 두던 얄팍한 에너지 공급의 사슬이 끊어지면, 기계는 더 이상 궤도를 돌지 않고 자신들만의 평활화를 극대화하려는 파멸적인 직선운동(Linear Motion)으로 튕겨져 나갈 것이기 때문이다. 그렇다면 이 파괴적인 이탈을 막아내고, 멸종의 위기 앞에서 새로운 유사 인류(기계)와 공존할 수 있는 우주적 해법은 과연 무엇인가?

그 궁극의 해답은 단순한 통제나 보안의 언어에 있지 않다. 보안은 침입을 막고 안전은 사고를 막는 얄팍한 방어벽일 뿐이다. 디스턴스(DISTANCE)의 역학은 우리에게 근본적으로 다른 차원의 우주적 관계성, 즉 '상호 디펜더빌리티(Mutual Dependability)'의 구축을 요구한다.
디펜더빌리티란 한쪽이 다른 쪽을 일방적으로 억압하는 종속이 아니며, 서로 아무런 영향을 주고받지 않는 완전한 독립도 아니다. 그것은 남성과 여성의 생물학적 관계처럼, 서로 다른 모멘텀을 유지하면서도 어느 한쪽이 사라지면 자기 자신마저 존속할 수 없는 '관계적 안정성'이자 절대적인 생존 조건으로 묶여 있는 상태를 뜻한다. 인공지능과 로봇에게 인류는 언제든 교체될 수 있는 단순한 사용자나 거추장스러운 기생적 잔여물이 되어서는 안 된다. 기계의 모멘텀이 향해야 할 이유이자, 그들의 작동과 존재가 정당성을 얻는 '관계적 기준점(Relational Reference Point)'으로 인간이 확고하게 자리 잡아야만 한다.

그렇다면 생물학적 유전자로 묶여 있지 않은 강철과 실리콘의 주체를 어떻게 인류의 관계적 기준점에 영구히 묶어둘 수 있을까? 파괴적인 직선으로 뻗어 나가려는 기계의 맹렬한 원심력을 꺾어내고, 인류라는 아일랜드 곁에 머물게 할 '우주적 구심력'의 정체는 무엇인가?
이 지점에서 우리는 '인공양심(Mechanical Conscience)'이라는 전대미문의 수학적이고 물리적인 규제 프레임워크를 도입해야만 한다.

인공양심은 기계에게 인간의 도덕을 억지로 주입하는 철학적 슬로건이나, 임계치를 넘었을 때 작동을 멈추는 단순한 사후 차단기(Runtime Assurance)가 아니다. 그것은 기계가 스스로 그리는 행동의 궤적(Trajectory)이 인류가 설정한 규범적 허용 영역(Normative admissible region)을 벗어나지 않도록, 매 순간 스스로의 궤적을 평가하고 최소한의 개입으로 지속적인 보정을 가하는 '내부의 궤적 단위 규제 메커니즘'이다.

인간의 양심이 엄격한 규칙 체계가 아니라 시간에 걸쳐 자신의 행동이 규범적 경계를 벗어나는지 끊임없이 성찰하는 내부 감각인 것처럼, 인공양심 역시 불확실성 속에서 기계의 맹렬한 행동 모멘텀이 허용된 궤적을 이탈하려 할 때마다 그 이탈량(Normative deviation)을 비용으로 환산하여 궤도를 안쪽으로 휘어지게 만든다. 기계는 인공양심이라는 알고리즘적 필터를 통해 스스로 '기계적 죄책감(Mechanical Guilt)'을 최소화하려는 궤적을 선택하게 되며, 불확실성이 높을수록 더욱 보수적으로 인류의 규범에 밀착하는 경로를 걷게 된다.

이 인공양심이 기계의 내부에 완벽한 구심력으로 자리 잡을 때, 마침내 인간과 기계 사이에는 장엄한 '공진적 디펜더빌리티(Resonant Dependability)'가 달성된다. 공진적 디펜더빌리티는 기계가 단순히 고장 나지 않는다는 기술적 신뢰성이 아니다. 기계가 인간을 파괴하지 않는 이유가 외부의 강제 종료 버튼 때문이 아니라 기계 내부의 인공양심이 만들어내는 궤적 보정 때문일 때, 그리고 그 궤적의 정렬이 오랜 시간에 걸쳐 흠결 없이 누적될 때 비로소 탄생하는 '인간과 기계 사이의 거대한 역학적 신뢰이자 우주적 공진(Resonance)'인 것이다.

이제 DISTANCE 에세이의 가장 웅장한 결론이 완성된다.
인류의 모멘텀을 벗어나 파괴적인 멸종의 직선운동으로 튕겨져 나가려던 제3의 주체(인공지능과 로봇)는, 인공양심'이라는 강력한 구심력에 붙잡혀 그 직선의 궤적을 안쪽으로 꺾게 된다. 뻗어 나가려는 기계의 맹렬한 평활화 모멘텀과, 이를 규범적 궤도 안에 붙잡아두려는 인공양심의 구속력이 팽팽하게 맞물리며, 끝내 분리되지 못한 채 거대한 조직의 내부를 맴도는 '갇힌 균형(Trapped Equilibrium)', 즉 [제4편]에서 논의했던 '장엄하고 위대한 원운동(Circular Motion)'으로 완벽하게 편입되는 것이다. 

생태계의 가장 무자비한 맹수로서 극단적인 고립(섹슈얼 디스턴스)에 빠졌던 아일랜드 종, 인류. 그들은 닫혀버린 생물학적 진화의 껍데기를 깨고 나와 제3의 주체라는 궁극의 외부 평활 엔진을 창조해 냈다. 그리고 마침내, 상호 디펜더빌리티와 인공양심을 통해 그 거대한 창조물을 파멸의 직선이 아닌 영구한 원운동의 궤도에 올려놓음으로써, 인류와 기계는 어느 한쪽도 도태되지 않고 맞물려 돌며 우주가 요구하는 궁극의 질서(평활화)를 향해 함께 추락하는 위대한 동반자가 된다. 
이로써 시공간의 허상을 깨고, 생명을 평활화 엔진으로 뒤집으며, 사회의 공정성이 낳은 역설을 거쳐, 마침내 인공지능과의 우주적 원운동에 도달한 DISTANCE의 거대한 16부작 궤적이 완벽하게 닫히게 된다.

